\clearpage
\begin{flushright}
  \fechaclase{10 de septiembre de 2026}
\end{flushright}

\subsection{Sliding Mode Control (SMC)}

Sea el sistema no lineal
\[
  \dot{X}_1=X_2
\]
\[
  \dot{X}_2=f(X_1,X_2,t)+v
\]
con la perturbación acotada
\[
  \lVert f(X_1,X_2,t)\rVert\leq L,
  \qquad L>0
\]
y las condiciones iniciales $X_{10}$ y $X_{20}$.

Se desea una dinámica estable, por ejemplo
\[
  \dot{X}_1+cX_1=0,
  \qquad c>0
\]
resolviendo
\[
  X_1=X_{10}\exp(-ct)
\]
y considerando que $\dot{X}_1=X_2$
\[
  X_2=-cX_{10}\exp(-ct)
\]
mostrando que esta dinámica hace estable al sistema. Es importante destacar
que esta dinámica no considera los efectos de la perturbación $r(X,t)$.

Para obtener esta dinámica definimos una nueva variable en el espacio de
estados llamada superficie de deslizamiento
\[
  \sigma=X_2+cX_1,
  \qquad c>0
\]
Note que
\[
  \dot{X}_1+c_1X_1=X_2+c_1X_1
\]
entonces si
\[
  \sigma=0
\]
el sistema sigue la dinámica deseada.

Se propone una función de Lyapunov que haga a la superficie asintóticamente
estable, para que $\sigma\to0$ conforme $t\to\infty$
\[
  V=\frac{1}{2}\sigma^2
\]
\[
  \dot{V}=\sigma\dot{\sigma}
\]

Derivando la superficie
\begin{align*}
  \dot{\sigma}
  &=\dot{X}_2+c\dot{X}_1\\
  &=f(X,t)+v+cX_2\\
  &=f(X,t)+v+cX_2\\
  &=f(X,t)+cX_2+v
\end{align*}
y se sustituye
\[
  \dot{V}=\sigma\left[f(X,t)+cX_2+v\right]
\]
\[
  \dot{V}=\sigma f(X,t)+\sigma[cX_2+v]
\]
\[
  \dot{V}\leq L\lvert\sigma\rvert+\sigma[cX_2+v]
\]

Haciendo
\[
  v=-cX_2+V,
\]
con el controlador virtual $V$
\[
  \dot{V}\leq L\lvert\sigma\rvert
  +\sigma(cX_2-cX_2+V)
\]
\[
  \dot{V}\leq L\lvert\sigma\rvert+\sigma V
\]

Elegimos
\[
  V=-K\operatorname{sign}(\sigma)
\]
donde $K$ es una ganancia a elegir y $\operatorname{sign}(\cdot)$ es
\[
  \operatorname{sign}(x)=
  \begin{cases}
    -1, & x<0,\\
    1, & x>0.
  \end{cases}
\]
También
\[
  \operatorname{sign}(\sigma)
  =\frac{\lvert\sigma\rvert}{\sigma}
  =\frac{\sigma}{\lvert\sigma\rvert}
\]
\[
  \dot{V}\leq L\lvert\sigma\rvert
  -K\sigma\frac{\lvert\sigma\rvert}{\sigma}
\]
\[
  \dot{V}\leq L\lvert\sigma\rvert-K\lvert\sigma\rvert
\]
\[
  \dot{V}\leq-(K-L)\lvert\sigma\rvert
\]
Si $K>L$ entonces
\[
  \dot{V}\leq-(K-L)\lvert\sigma\rvert<0
\]
haciendo al sistema globalmente asintóticamente estable con el controlador
\[
  \sigma=X_2+cX_1
\]
\[
  v=-cX_2-K\operatorname{sign}(\sigma)
\]